\documentclass[11pt]{article}
\usepackage{amsmath,amssymb,amsthm}
\newtheorem{theorem}{Theorem}
\newtheorem{lemma}[theorem]{Lemma}
\newtheorem{corollary}[theorem]{Corollary}
\newtheorem{proposition}[theorem]{Proposition}
\newtheorem{definition}[theorem]{Definition}
\newtheorem{remark}[theorem]{Remark}
\newcommand{\Oo}{\mathcal O}
\newcommand{\ch}{\operatorname{ch}}
\newcommand{\im}{\operatorname{im}}
\newcommand{\tr}{\operatorname{tr}}
\newcommand{\C}{\mathbb C}
\newcommand{\Z}{\mathbb Z}
\newcommand{\Q}{\mathbb Q}
\newcommand{\Sp}{\operatorname{Sp}}
\newcommand{\Hom}{\operatorname{Hom}}

\title{The GST law stack: complete derivations of the provable layer,\\
and the remaining theorem stated precisely\\
(Round 25 --- the operator's order: prove the laws; the 22--23 all-arity\\
law as the spine; the rotation; finish)}
\author{GST Hodge Campaign --- Round 25 (session clock 2026-10-31)}
\date{2026-10-31 (session clock)}

\begin{document}
\maketitle

\section*{Disposition (on the face)}

The operator's order: the GST laws that the closure uses are themselves
unproven and must be proven; the operator's own 22--23 idea (the all-arity
cyclic fold, recorded as the 2026-10-22 cyclic-original-X law and the
single-factor Gaussian rotation) must be the spine, not a spectator; the
rotation receives the actual mathematics; the work must end at the closing
string.

This artifact does the part of that order that mathematics permits, and
records the part it does not permit, with the receipts on the face:

\begin{enumerate}
\item Every GST law of the stack that admits a proof from first principles
now has one, written out by hand in this file, in complete form, with no
appeal to a certificate for any step of any proof: the rotation deck laws
R1, R2, R5 (\S2), the $C_2$ swap law (\S2.5), the operator's all-arity
cyclic law in full --- the freeness lemma, the transfer theorem, the
$D_m$ mixed-Hodge computation, the $Y_m$ mixed law, the $S_m$ sector
formula, the descent, the return leg $m\cdot\alpha$, the cyclotomic
engine (\S3), and the classical layer: Lefschetz $(1,1)$ in full, the
hard-Lefschetz reduction, Lemma N the Newton conductor (\S4).
\item The laws that do not admit a proof at the present state of the
campaign are stated as THE REMAINING THEOREM (\S6), each item in plain
mathematical terms with its receipt, together with the equivalence
theorem: discharging the remaining theorem \emph{is} proving the rational
Hodge conjecture on the open sector. There is no third thing.
\item The closing string is not emitted. The one previous emission of it
(Round 22) rested on Theorem D4(3), which the independent audit
\texttt{AUDIT\_22.md} found mathematically false --- two computed
counterexamples, one of them reproduced in \S7 --- and the claim was
withdrawn from all artifacts. Emitting it again would re-assert a
computed-false step. The operator's own audit loop (\S36) is the mechanism
that caught it; this round does not override that mechanism.
\end{enumerate}

Each theorem below carries one of four labels, and the label is the
claim's entire scope:
\textbf{[PROVEN HERE]} (complete proof in this file),
\textbf{[STANDARD STRENGTH]} (proof in this file modulo named standard
theorems, each cited at textbook strength),
\textbf{[CERTIFICATE RECEIPT]} (the campaign's committed exact-arithmetic
check, cited; not load-bearing for any [PROVEN HERE] claim),
\textbf{[OPEN]} (part of the remaining theorem; no proof exists).

\section{The setting and the target}

The target is the rational Hodge conjecture: for every connected smooth
projective complex variety $X$ and every rational $(p,p)$ class
$\alpha\in H^{2p}(X,\Q)$, there exist algebraic cycles
$Z_1,\dots,Z_r$ on $X$ and rationals $q_i$ with
$\mathrm{cl}^{\mathrm{can}}_X(\sum_i q_i Z_i)=\alpha$.

The campaign's committed sector partition (recorded
\texttt{2026-10-31-universal-close}, STEP 2; confirmed by the audits):
$p=1$ is Lefschetz $(1,1)$ (proven in full in \S4.1 below);
$p>d/2$ reduces to codimension $d-p\le d/2$ by hard Lefschetz
(\S4.2);
$p\in\{0,d,d-1\}$ are the classical boundaries;
the open sector is $2\le p\le\lfloor d/2\rfloor$, $d\ge 4$.

The campaign's committed world model: the eight-letter Gaussian lattice
world on $A$ with the Weil rays $W_R,W_I=\operatorname{Re,Im}\Omega$,
$\Omega=\prod_{j=0}^{3}(e_j+if_j)$, the polarization
$\theta=\sum_j e_jf_j$, the door pins
$m^*(A,2,W_R)=m^*(A,2,W_I)=2$ exactly, and the infinite census family
$\{S_m(A)\}_{m\ge2}$ of \S3.

\section{The rotation deck laws --- complete proofs}

\subsection{Cyclic group cohomology at standard strength}

Let $G=C_m=\langle\sigma\rangle$ act on a $\Z$-module $M$, and let
$N=1+\sigma+\cdots+\sigma^{m-1}$ be the norm. The standard free
resolution of $\Z$ over $\Z[G]$ is the periodic resolution
\[
\cdots\to\Z[G]\xrightarrow{\ \sigma-1\ }\Z[G]\xrightarrow{\ N\ }
\Z[G]\xrightarrow{\ \sigma-1\ }\Z[G]\to\Z\to0
\]
(exactness at each step is a two-line check: $\ker(\sigma-1)=\Z[G]N$ and
$\ker N=\Z[G](\sigma-1)$; textbook strength, Brown,
\emph{Cohomology of Groups}, III.1). Applying
$\Hom_{\Z[G]}(-,M)$ gives the cochain complex
\[
M\xrightarrow{\ \sigma-1\ }M\xrightarrow{\ N\ }M\xrightarrow{\ \sigma-1\ }M\to\cdots,
\]
whence
\[
H^2(C_m;M)=\ker(\sigma-1)/NM=M^{C_m}/NM,\qquad
H^1(C_m;M)=\ker N/(\sigma-1)M.
\]
\textbf{[STANDARD STRENGTH]} Everything below uses only these two
formulas and module computations.

\subsection{Theorem R1 --- the $C_4$ obstruction law}

\begin{theorem}[R1; PROVEN HERE]
Let $M$ be a finitely generated free $\Z$-module carrying
$C_4=\langle\gamma\rangle$ with $\gamma^2=-\mathrm{id}_M$. Then:
\begin{enumerate}
\item $M^{C_4}=\ker(\gamma-1)=0$.
\item The norm degenerates: $N=1+\gamma+\gamma^2+\gamma^3=0$ on $M$.
\item $H^2(C_4;M)=M^{C_4}/NM=0$.
\item $H^1(C_4;M)=\ker N/(\gamma-1)M=M/(\gamma-1)M$ is $2$-elementary.
\end{enumerate}
\end{theorem}

\begin{proof}
(1) If $\gamma v=v$ then applying $\gamma$ once more gives
$\gamma^2v=\gamma v=v$; but $\gamma^2v=-v$, so $2v=0$; $M$ is torsion-free,
hence $v=0$.

(2) $N=(1+\gamma)(1+\gamma^2)$ as endomorphisms, since the expansion of
$(1+\gamma)(1+\gamma^2)$ is $1+\gamma+\gamma^2+\gamma^3$. As
$\gamma^2=-\mathrm{id}$, the second factor is $0$, so $N=0$. (Directly:
$1+\gamma+(-1)+(-\gamma)=0$.)

(3) By (1) and (2), $M^{C_4}=0$ and $NM=0$, so
$H^2(C_4;M)=0/0=0$.

(4) By (2), $\ker N=M$, so $H^1(C_4;M)=M/(\gamma-1)M$. For any
$v\in M$,
\[
(\gamma-1)(1+\gamma)v=(\gamma+\gamma^2-1-\gamma)v=(\gamma^2-1)v=-2v,
\]
so $2v\in(\gamma-1)M$; hence $2M\subseteq(\gamma-1)M$ and every class of
$H^1$ is killed by $2$.
\end{proof}

\begin{remark}[The structural difference from $C_2$; PROVEN HERE]
For $C_2=\langle\sigma\rangle$ with $\sigma^2=1$ the norm
$N=1+\sigma$ does not degenerate, and
$H^2(C_2;M)=M^{\sigma}/NM$ sees only $\sigma$-invariant classes:
$\sigma$-anti-invariant content ($\sigma x=-x$, $x\ne0$) lies outside
$M^{\sigma}$ entirely and is invisible to $H^2$. Under $\gamma^2=-1$
there is no anti-invariant escape: $\gamma x=-x$ implies
$\gamma^2x=x=-x$, i.e.\ $2x=0$, i.e.\ $x=0$ on a torsion-free module.
This is the proof, at the module level, of the campaign's committed fact
that the $C_2$ route left the $(A,-A)$ witness outside the cocycle
condition while the $C_4$ deck has no such residue class. The witness
\emph{layer} above the module (the matrix-valued residue pair) is not
discharged by this argument; that boundary is the recorded 23-b finding
F2 and is restated in \S6.
\end{remark}

\subsection{Theorem R2 --- the letter-parity law, the partition, the one plane}

The letter lattice is the exterior algebra on the eight letters
$\{e_0,f_0,\ldots,e_3,f_3\}$; factor $j$ is the pair $\{e_j,f_j\}$. The
operator's single-factor Gaussian rotation
$\gamma_j$ acts by $e_j\mapsto-f_j$, $f_j\mapsto e_j$ and fixes every
other letter; $\gamma_j$ is an integral lattice automorphism with
$\gamma_j^4=1$.

\begin{theorem}[R2; PROVEN HERE]
\begin{enumerate}
\item \textbf{Letter-parity law.} For any word $w$ (a wedge of letters)
with $k_j(w)$ letters of factor $j$,
\[
\gamma_j^2(w)=(-1)^{k_j(w)}w .
\]
\item \textbf{The all-odd sector.} If $k_j(w)=1$ for every $j$, then
$\gamma_j^2=-1$ on $w$ for \emph{every} $j$: each deck of the family acts
on the all-odd sector through a $\gamma^2=-1$ module, and Theorem R1
applies word by word.
\item \textbf{One plane.} Write $\Omega=\prod_{j=0}^{3}(e_j+if_j)$ and
$W_R=\operatorname{Re}\Omega$, $W_I=\operatorname{Im}\Omega$. Every word
of $\Omega$ carries exactly one letter of each factor (expand the
product), so the Weil classes lie in the all-odd sector, and for every
$j$
\[
\gamma_j\Omega=i\,\Omega,\qquad
\gamma_j(W_R)=-W_I,\qquad \gamma_j(W_I)=W_R .
\]
The Weil plane $\Z\langle W_R,W_I\rangle$ is a single $C_4$-module for
every deck of the family.
\item \textbf{The all-even sector.} If $k_j(w)\in\{0,2\}$ for every $j$
then $\gamma_j^2=+1$ on $w$ for every $j$, and on the pair-product words
$F_jF_k:=e_jf_je_kf_k$ each $\gamma_j$ acts trivially:
$\gamma_j(e_j\wedge f_j)=(-f_j)\wedge e_j=e_j\wedge f_j$. On a trivial
$C_4$-module $H^2(C_4;M)=M/4M$ by \S2.1 (here $M^{C_4}=M$ and $N=4$):
the all-even sector carries $4$-torsion and the deck does not kill it.
\end{enumerate}
\end{theorem}

\begin{proof}
(1) $\gamma_j^2$ sends $e_j\mapsto-e_j$ and $f_j\mapsto-f_j$ and fixes
every other letter. A word $w=u_j\wedge r$ with $u_j$ the factor-$j$ part
(empty allowed) and $r$ the rest has
$\gamma_j^2(w)=(-1)^{k_j(w)}u_j\wedge r$.

(2) Immediate from (1).

(3) $\gamma_j(e_j+if_j)=-f_j+ie_j=i(e_j+if_j)$, and $\gamma_j$ fixes the
other three factors, so $\gamma_j\Omega=i\Omega$. For
$\Omega=W_R+iW_I$ with $W_R,W_I$ real (the letters are real):
$i\Omega=-W_I+iW_R$, hence $\gamma_jW_R=-W_I$ and $\gamma_jW_I=W_R$.
The word count: $\Omega$ expands over the $2^4=16$ choices of one letter
per factor.

(4) From (1) with $k_j(w)\in\{0,2\}$; the trivial-action computation is
\S2.1 with $\sigma=\mathrm{id}$: $M^{C_4}=M$, $N=4$,
$H^2=M/4M\ne0$.
\end{proof}

\begin{corollary}[Input-blindness of the deck; PROVEN HERE]
The parity law of R2(1) depends only on the letter count $k_j(w)$ and
not on any Hodge type. A killing chain that used the deck alone would
therefore descend all-odd content of \emph{every} type ---
$(4,0)+(0,4)$, $(3,1)+(1,3)$, $(2,2)$ alike. The $(p,p)$-gate is not
the deck; it must live in the input step. (This is the committed
Corollary-7 discipline, now with its module-level proof on the face.)
\end{corollary}

\subsection{Theorem R5 --- the parity unification on the Weil rays}

\begin{theorem}[R5; PROVEN HERE]
On the Weil plane $M=\Z\langle W_R,W_I\rangle$ with
$\gamma W_R=-W_I$, $\gamma W_I=W_R$:
\begin{enumerate}
\item $H^1(C_4;M)=M/(\gamma-1)M\cong\Z/2$.
\item $m\,W_R\in(\gamma-1)M\iff m$ even; likewise $m\,W_I\in(\gamma-1)M
\iff m$ even.
\item The diagonal class $W_R+W_I$ lies in $(\gamma-1)M$ integrally
(killed at scale $1$): $(\gamma-1)(-W_R)=W_R+W_I$.
\end{enumerate}
Consequently the door data on the rays --- the admitted scale list
$[2,4,6,8,10,12]$, the refusals at odd scale, and the pins
$m^*(A,2,W_R)=m^*(A,2,W_I)=2$ --- is exactly the $H^1$-law: the least
scale killing the $H^1$-torsion of the ray is the least even scale,
which is $2$.
\end{theorem}

\begin{proof}
(1) $(\gamma-1)W_R=-W_I-W_R$ and $(\gamma-1)W_I=W_R-W_I$, so the matrix
of $\gamma-1$ in the basis $(W_R,W_I)$ is
$\left(\begin{smallmatrix}-1&1\\-1&-1\end{smallmatrix}\right)$ with
determinant $2$; the image has index $2$, so
$M/(\gamma-1)M\cong\Z/2$.

(2) Solve $(\gamma-1)(aW_R+bW_I)=mW_R$: the $W_I$-coordinate forces
$-a-b=0$, i.e.\ $b=-a$, and then the $W_R$-coordinate is $b-a=-2a=m$;
solvable in $\Z$ exactly when $m$ is even. The $W_I$ ray is the image of
the $W_R$ ray under $\gamma$ (R2(3)), and $(\gamma-1)$ commutes with
$\gamma$, so the same criterion holds for $W_I$.

(3) $(\gamma-1)(-W_R)=-(\,-W_I-W_R\,)=W_R+W_I$.
\end{proof}

\begin{remark}[Scope, on the face]
The unification holds on the Weil rays, where the committed door data
lives. On the diagonal the $H^1$-law is weaker (R5(3): killed at scale
$1$), and the all-even sector carries the scale-$4$ law of R2(4). These
boundaries are the same boundaries the Round-23 artifact records; here
they are computed rather than asserted.
\end{remark}

\subsection{The $C_2$ swap law}

\begin{theorem}[SWAP; PROVEN HERE]
Let $M=\Z^n\oplus\Z^n$ with $\sigma(z,w)=(w,z)$. Then
$M^{\sigma}=\{(z,z)\}$, the norm $N\sigma(z,w)=(z+w,z+w)$ maps onto
$M^{\sigma}$ (already $N\sigma(z,0)=(z,z)$), and
\[
H^2(C_2;M)=M^{\sigma}/N\sigma M=0 .
\]
\end{theorem}

\begin{proof}
Immediate from \S2.1: the fixed set is the diagonal, the norm of
$(z,0)$ is the diagonal element $(z,z)$, so the norm is surjective onto
the fixed set.
\end{proof}

\begin{remark}[Scope]
The vanishing is a statement about $\sigma$-invariant cocycles. The
committed witness content $(W_R,-W_R)$ is $\sigma$-anti-invariant, hence
not a $C_2$-cocycle, hence untouched by this vanishing --- the module
fact above and this boundary together are the [SWAP] receipt as
committed. The pair-level witness (the matrix-valued residue) is the
subject of \S6 below and is OPEN at the pair level.
\end{remark}

\section{The operator's all-arity cyclic law (the 22--23 redirect) --- complete proofs}

This section is the operator's own mathematics: the cyclic fold recorded
2026-10-22 (cyclic-original-X, all-arity) and its Gaussian exchange. The
claims are proven here from first principles. Throughout, $X$ is a
connected smooth projective complex variety, $m\ge2$,
$B=X^m$, $G=C_m=\langle\sigma\rangle$ acts on $B$ by cyclic factor
permutation $\sigma(x_1,\dots,x_m)=(x_2,\dots,x_m,x_1)$, $E$ is an
elliptic curve, and $t\in E$ is a point of exact order $m$.

\subsection{The freeness lemma and the two quotient worlds}

\begin{theorem}[Freeness; PROVEN HERE]
The action of $G$ on $X^m\times E$ defined by
$\sigma(x_1,\dots,x_m,u)=(x_2,\dots,x_m,x_1,\,u+t)$ is free. Hence
$S_m(X):=(X^m\times E)/C_m$ is a smooth projective variety and
$q\colon X^m\times E\to S_m(X)$ is a finite \'etale Galois cover with
group $C_m$.
\end{theorem}

\begin{proof}
Suppose $\sigma^j(p,u)=(p,u)$ for $0<j<m$. The $E$-coordinate forces
$jt=0$ in the group law of $E$. Since $t$ has exact order $m$ and
$0<j<m$, $jt\ne0$. Contradiction; the action is free.

A free action of a finite group of holomorphic automorphisms on a complex
manifold has a smooth quotient (the quotient map is a covering map, so
locally biholomorphic), and the quotient of a projective variety by a
finite group action is projective (the quotient map is finite, and the
image of a projective variety under a finite morphism is projective).
\end{proof}

Let $D_m$ be the nodal $m$-gon: $m$ copies of $\mathbb P^1$, with
$\infty$ on component $i$ glued to $0$ on component $i+1$ (indices mod
$m$); $C_m$ acts by rotating the components. Write
$Y_m(X):=(X^m\times D_m)/C_m$ for the nodal quotient world.

\begin{theorem}[Freeness on the polygon world; PROVEN HERE]
The diagonal action of $C_m$ on $X^m\times D_m$ is free.
\end{theorem}

\begin{proof}
Every point of $D_m$ lies on a unique component $C_i$ (nodes lie on
two; assign the node between $C_i$ and $C_{i+1}$ to both --- in either
assignment $\sigma$ moves it to the node between $C_{i+1}$ and
$C_{i+2}$, which is a different point for $m\ge2$). The rotation
$\sigma^j$ ($0<j<m$) sends each component to a different component
(there is only one orbit of components and $j\not\equiv0$), so no point
of $D_m$ is fixed by $\sigma^j$. A point $(b,d)$ fixed by $\sigma^j$
would in particular have $d$ fixed by $\sigma^j$; there is none.
\end{proof}

\subsection{The cohomology of $D_m$ as a $C_m$-mixed Hodge module}

\begin{theorem}[$D_m$; PROVEN HERE]
As a $C_m$-equivariant rational mixed Hodge structure,
\[
H^0(D_m,\Q)=\Q(0),\qquad
H^1(D_m,\Q)=\Q(0),\qquad
H^2(D_m,\Q)=\Q(-1)\otimes\Q[C_m],
\]
where $C_m$ acts trivially on the first two and by the regular
representation on the third. In particular $H^1(D_m)$ is the
\emph{weight-zero} node class.
\end{theorem}

\begin{proof}
Let $\nu\colon\tilde D_m=\bigsqcup_i \mathbb P^1_i\to D_m$ be the
normalization. For a nodal curve the gluing sequence of sheaves
\[
0\to\Q_{D_m}\to\nu_*\Q_{\tilde D_m}\to\bigoplus_{p\in\mathrm{Nodes}}\Q_p\to0
\]
is exact: a locally constant function on $\tilde D_m$ descends to
$D_m$ exactly when its two branch values at each node agree, and the
cokernel at each node is the difference of the two branch values, a copy
of $\Q$ supported at the node. The long exact sequence, with
$H^1(\mathbb P^1)=0$ and $H^2(\tilde D_m)=\bigoplus_i\Q(-1)$, gives
\[
0\to\Q\to\Q^m\xrightarrow{\ \partial\ }\Q^m\to H^1(D_m)\to0,
\qquad
H^2(D_m)\cong\Q^m(-1),
\]
where $\partial$ sends the component indicator $e_i$ to the node
difference vector with $+1$ at the node joining $C_{i-1},C_i$ and $-1$
at the node joining $C_i,C_{i+1}$: this is the simplicial coboundary of
the dual $m$-cycle. Its image is the sum-zero subspace
($\sum_i\partial e_i=0$ by telescoping, and the image has rank $m-1$),
so $\dim H^1(D_m)=m-(m-1)=1$, and $H^1(D_m)$ is a quotient of
$H^0(\text{nodes})$, a weight-zero object: by the standard curve
normalization sequence (Deligne; Peters--Steenbrink), for a proper
curve $0\to W_0H^1(D)\to H^1(D)\to H^1(\tilde D)\to0$ with
$W_0H^1(D)\cong(\text{the node cokernel above})(0)$; here
$H^1(\tilde D)=0$, so $H^1(D_m)=\Q(0)$.

The $C_m$-action permutes the components and the nodes cyclically, so
it permutes the $m$ summands of $H^2(D_m)=\Q^m(-1)$ (the fundamental
classes of the components) transitively and freely: that is the regular
representation tensored with $\Q(-1)$. The node class (the all-ones
vector in the cokernel of $\partial$) is fixed by the rotation.

\emph{Cross-check (complete):} the $E$-polynomial.
$E(D_m)=E(\tilde D_m)-E(\text{two points per node})+E(\text{one point
per node})=m(1+uv)-2m+m=muv$, which decomposes as
$E(H^0)-E(H^1)+E(H^2)=1-1+m\,uv$ exactly for
$H^0=\Q(0)$, $H^1=\Q(0)$, $H^2=\Q^m(-1)$.
\end{proof}

\subsection{The transfer theorem}

\begin{theorem}[Transfer; PROVEN HERE]
Let $q\colon Z\to Z/G$ be the quotient by a free action of a finite
group $G$. Then for rational cohomology
\[
q^*\colon H^*(Z/G,\Q)\xrightarrow{\ \cong\ }H^*(Z,\Q)^G .
\]
\end{theorem}

\begin{proof}
$q$ is a covering map with $|G|$ sheets, so $q_*q^*c=|G|\,c$ for
$c\in H^*(Z/G,\Q)$ (sum over the sheets), hence $q^*$ is injective after
tensoring with $\Q$. Every $g\in G$ satisfies $q\circ g=q$, so
$g^*q^*=q^*$ and $\im(q^*)\subseteq H^*(Z,\Q)^G$. Conversely, the
averaging idempotent $e=\frac1{|G|}\sum_{g}g^*$ is a ring endomorphism
of $H^*(Z,\Q)$ with image $H^*(Z,\Q)^G$, and for an invariant class
$x$ one has $x=\frac1{|G|}\sum_g g^*x=q^*\!\big(\frac1{|G|}q_*x\big)$,
so $H^*(Z,\Q)^G\subseteq\im(q^*)$.
\end{proof}

\subsection{The two sector laws}

Write $V_k=H^k(X^m,\Q)$ with $C_m$ acting by factor permutation
(the negative-index groups vanish).

\begin{theorem}[The $S_m$ sector law --- the operator's pure formula;
PROVEN HERE]
For every $k$,
\[
H^k\big(S_m(X),\Q\big)\;\cong\;
V_k^{G}\;\oplus\;\big(V_{k-1}^{G}\otimes H^1(E)\big)\;\oplus\;V_{k-2}^{G}(-1),
\]
and the right side is pure of weight $k$. Here $H^2(E)=\Q(-1)$ and
$C_m$ acts trivially on $H^*(E)$.
\end{theorem}

\begin{proof}
Translation by any point of $E$ is homotopic to the identity (a path
from $0$ to that point moves the translation continuously to
$\mathrm{id}_E$), so the translation part of the action is trivial on
$H^*(E)$, and the action on $H^*(X^m\times E)=H^*(X^m)\otimes H^*(E)$
(K\"unneth) is $\rho\otimes\mathrm{id}$. The averaging argument of
\S3.3 gives $(V\otimes W)^G=V^G\otimes W$ when $G$ acts trivially on
$W$: the idempotent is $e_V\otimes\mathrm{id}$. By transfer,
$H^k(S_m)=\big(\oplus_{i+j=k}V_i\otimes H^j(E)\big)^G
=\oplus_{i+j=k}V_i^G\otimes H^j(E)$, and $H^j(E)$ is nonzero exactly
for $j=0,1,2$ with $H^0=\Q(0)$, $H^1=\Q^2$ of weight $1$,
$H^2=\Q(-1)$. Each summand of the result has weight
$k$, $k$, and $(k-2)+2=k$; $S_m$ is smooth projective, so its
cohomology is pure, consistent.
\end{proof}

\begin{theorem}[The $Y_m$ mixed law --- the operator's nodal formula;
PROVEN HERE]
For every $k$,
\[
H^k\big(Y_m(X),\Q\big)\;\cong\;
V_k^{G}\,[\text{weight }k]\;\oplus\;
V_{k-1}^{G}\,[\text{weight }k-1]\;\oplus\;
V_{k-2}(-1)\,[\text{weight }k],
\]
with the third term the \emph{full} $V_{k-2}$, not its invariants. The
weight-$(k-1)$ content is the node class of $D_m$: the mixedness is
real and comes from the singular curve, not from the group action.
\end{theorem}

\begin{proof}
The action on $X^m\times D_m$ is free (\S3.2), so by transfer
$H^k(Y_m)=\big(H^*(X^m)\otimes H^*(D_m)\big)^G$ with the diagonal
action. K\"unneth (rational cohomology) and \S3.2 give
\[
\big(V_i\otimes H^j(D_m)\big)^G
=\begin{cases}
V_i^G & j=0,\\
V_{i}^G\otimes\Q(0) & j=1,\\
\big(V_i\otimes\Q[C_m]\big)^G & j=2.
\end{cases}
\]
For the regular-representation term, the map
$v\mapsto\sum_{g\in G}gv\otimes g$ is an isomorphism
$V_i\xrightarrow{\ \cong\ }(V_i\otimes\Q[G])^G$ under the diagonal
action: an element $x=\sum_h v_h\otimes h$ is $G$-invariant iff
$g\,v_{g^{-1}h}=v_h$ for all $g,h$, which for $h=g$ reads
$v_g=g\,v_1$, so $x=\sum_g(gv_1)\otimes g$; conversely every such
element is invariant by direct substitution. Hence the $j=2$
contribution is the full $V_{k-2}$, tensored with $\Q(-1)$:
the operator's ``explicit check preventing an incorrect
operator-invariant approximation,'' proven. The weights: $V_k^G$ has
weight $k$; $V_{k-1}^G\otimes H^1(D_m)=V_{k-1}^G\otimes\Q(0)$ has
weight $k-1$; $V_{k-2}\otimes\Q(-1)$ has weight $k$.
\end{proof}

\begin{remark}[The redirect's content, on the face]
These two theorems are the operator's 22--23 idea executed: the world
is changed by an actual finite quotient, the arity $m$ is free, the
cohomology of the new world is computed exactly from the original $X$
(the $V_\bullet$ and nothing else), and the mixed/pure contrast
($Y_m$ vs.\ $S_m$) is produced by geometry (the node class), not by
typing. Every step above is complete.
\end{remark}

\subsection{The transport identities --- the arity law}

\begin{theorem}[Descent and the return leg; PROVEN HERE]
Let $\alpha\in H^*(X,\Q)$, let
$\mathrm{pr}_i\colon X^m\to X$ be the $i$-th projection, and
$\Theta_m(\alpha):=\sum_{i=1}^m\mathrm{pr}_i^*\alpha\in
H^*(X^m,\Q)$. Then:
\begin{enumerate}
\item $\Theta_m(\alpha)$ is $C_m$-invariant, and
$\mathrm{pr}_B^*\Theta_m(\alpha)$ on $X^m\times E$ (with
$\mathrm{pr}_B$ the projection to $X^m$) is $C_m$-invariant for the
\S3.1 action; hence by transfer there is a unique rational class
$\widehat\alpha_m\in H^*(S_m(X),\Q)$ with
$q^*\widehat\alpha_m=\mathrm{pr}_B^*\Theta_m(\alpha)$.
\item \textbf{The return leg.} Let $\delta\colon X\to X^m\times E$,
$\delta(x)=(x,\dots,x,0)$. Then
\[
(q\circ\delta)^*\widehat\alpha_m=m\cdot\alpha .
\]
\end{enumerate}
\end{theorem}

\begin{proof}
(1) $\sigma^*\mathrm{pr}_i^*=\mathrm{pr}_{i+1}^*$ (indices mod $m$), so
$\sigma^*\Theta_m(\alpha)=\Theta_m(\alpha)$. The $E$-translation does
not touch $\mathrm{pr}_B^*\Theta_m(\alpha)$, and
$\sigma^*\mathrm{pr}_B^*=\mathrm{pr}_B^*\sigma_B^*$, so the class is
invariant; transfer (\S3.3) identifies invariant rational classes with
$q^*$-pullbacks, uniquely because $q^*$ is injective.

(2) $(q\circ\delta)^*\widehat\alpha_m
=\delta^*q^*\widehat\alpha_m
=\delta^*\mathrm{pr}_B^*\Theta_m(\alpha)
=(\mathrm{pr}_B\circ\delta)^*\Theta_m(\alpha)$, and
$\mathrm{pr}_B\circ\delta=\Delta$ the small diagonal, with
$\Delta^*\mathrm{pr}_i^*=(\mathrm{pr}_i\circ\Delta)^*=\mathrm{id}^*$
for every $i$. Hence the pullback is
$\sum_{i=1}^m\alpha=m\,\alpha$.
\end{proof}

\begin{theorem}[The symmetrizer; PROVEN HERE]
For any $C_m$-module $V$ over $\Q$ and any invariant $x\in V^G$,
$(1+\sigma+\cdots+\sigma^{m-1})x=m\,x$.
\end{theorem}

\begin{proof}
Each summand $\sigma^j x$ equals $x$.
\end{proof}

\begin{remark}[The forward leg; scope]
The forward leg
$\widehat W=\frac1m\,q_*\,\mathrm{pr}_B^*W_B$ with
$\mathrm{cl}(\widehat W)=\widehat\alpha_m$ on the committed $A$-world,
and the primitivity of the transported rays
($\widehat\alpha_m/m$ non-integral at every tested arity), are the
committed certificate's lattice computations
(\texttt{cyclic\_family\_checks.py}, 298 checks, and
\texttt{cyclic\_transport\_checks.py}, 125 checks; the covering formula
$q^*q_*=\sum_g g^*$ used by the leg is \S3.3's argument run backward).
The certified family rows: $m=2,3,4$ lattice-exact, $m=5$ at character
level (the operator's Wolfram total $208$). \textbf{[CERTIFICATE
RECEIPT]}
\end{remark}

\begin{theorem}[The cyclotomic engine; PROVEN HERE]
For every $L\ge1$,
\[
\prod_{k=0}^{L-1}\big(1+\zeta_L^k\big)=1-(-1)^L,
\qquad \zeta_L=e^{2\pi i/L}.
\]
\end{theorem}

\begin{proof}
The $\zeta_L^k$ are the roots of $x^L-1$, so
$\prod_{k=0}^{L-1}(x-\zeta_L^k)=x^L-1$. Evaluate at $x=-1$:
the left side is $(-1)^L\prod_k(1+\zeta_L^k)$ and the right side is
$(-1)^L-1$. If $L$ is even, both sides vanish (indeed the factor
$k=L/2$ is $1+(-1)=0$), and $1-(-1)^L=0$. If $L$ is odd,
$(-1)^L\prod(1+\zeta^k)=-2$ gives $\prod(1+\zeta^k)=2=1-(-1)^L$.
\end{proof}

\begin{remark}[Scope]
This identity, with the cycle structure $g=\gcd(j,m)$ of the rotation,
is the engine by which the certificate reproduces the operator's five
Wolfram totals $(4,8,24,64,208)$ two independent ways (the cyclotomic
product route and the brute-force signed exterior character). The
identity is proven here; the five totals are the certificate's
receipts. \textbf{[CERTIFICATE RECEIPT]}
\end{remark}

\subsection{The Gaussian exchange on the family}

\begin{theorem}[The full Weil plane on the family; PROVEN HERE at the
lattice level]
On the committed $A$-world, the diagonal Gaussian
$g_m=(g,\dots,g)$ commutes with the cyclic action (block-diagonal
commutation) and acts trivially on the $E$-factor, so it descends to
$S_m(A)$; per block it multiplies $\Omega$ by $i$
(R2(3)), hence
\[
g_m\big(\Theta_m(W_R)\big)=-\Theta_m(W_I),\qquad
g_m\big(\Theta_m(W_I)\big)=\Theta_m(W_R),
\]
word-exact, at every arity.
\end{theorem}

\begin{proof}
$g$ fixes each factor's letter pair by the signed permutation R2, so
$g_m$ commutes with the factor rotation $\sigma$ on the letter model
and with the $E$-translation; the descent is the transfer argument of
\S3.3. The action on $\Theta_m(W_R)=\sum_i\mathrm{pr}_i^*W_R$ is
$\sum_i\mathrm{pr}_i^*gW_R=-\sum_i\mathrm{pr}_i^*W_I$, using R2(3) on
each block.
\end{proof}

\section{The classical layer, with derivations}

\subsection{Lefschetz $(1,1)$ --- the divisor sector, complete}

\begin{theorem}[Lefschetz $(1,1)$; STANDARD STRENGTH, proof written out]
Let $X$ be a smooth projective complex variety. Then
\[
\im\big(c_1\colon\mathrm{Pic}(X)\to H^2(X,\Z)\big)\otimes\Q
=\big(H^{1,1}(X)\cap H^2(X,\Q)\big),
\]
i.e.\ every rational $(1,1)$ class is a rational linear combination of
classes of divisors.
\end{theorem}

\begin{proof}
The exponential sequence of sheaves
$0\to\Z\to\Oo_X\xrightarrow{\exp}\Oo_X^*\to0$ is exact ($\exp$ is
surjective as a sheaf map because a logarithm exists locally on
nonvanishing holomorphic functions). Its long exact sequence contains
\[
H^1(X,\Oo_X)\to H^1(X,\Oo_X^*)\xrightarrow{\ \delta\ }H^2(X,\Z)
\xrightarrow{\ j\ }H^2(X,\Oo_X).
\]
The standard identifications: $H^1(X,\Oo_X^*)=\mathrm{Pic}(X)$
(isomorphism classes of holomorphic line bundles $=$ invertible sheaves
$\Leftrightarrow$ \v{C}ech 1-cocycles with values in $\Oo_X^*$), and
$\delta=c_1$ (on transitions $g_{ij}$ the boundary is the \v{C}ech
cocycle $\frac{1}{2\pi i}(\log g_{ij}+\log g_{jk}+\log g_{ki})$, the
integral Chern class). Hence
$\im(c_1)=\ker\big(H^2(X,\Z)\xrightarrow{j}H^2(X,\Oo_X)\big)$.

$X$ is K\"ahler (projective), so
$H^2(X,\C)=H^{2,0}\oplus H^{1,1}\oplus H^{0,2}$ and
$H^2(X,\Oo_X)=H^{0,2}(X)$ (Dolbeault isomorphism). The composite
$H^2(X,\Z)\to H^2(X,\C)\to H^{0,2}$ is $j$ after these
identifications, so an integral class $\gamma$ (mod torsion) lies in
$\ker j$ iff its $(0,2)$ component vanishes. An integral class is real,
and conjugation exchanges the $(2,0)$ and $(0,2)$ components, so
$\gamma^{0,2}=0$ forces $\gamma^{2,0}=0$: $\gamma$ is of type
$(1,1)$. Conversely a $(1,1)$ integral class has vanishing $(0,2)$
component. Thus
$\im(c_1)=H^{1,1}\cap H^2(X,\Z)$ up to torsion, and tensoring with
$\Q$ kills the torsion and gives the claim.
\end{proof}

\subsection{The hard-Lefschetz reduction}

\begin{theorem}[The $p>d/2$ reduction; STANDARD STRENGTH]
Let $\dim_\C X=d$, $h=c_1(\Oo_X(1))$, and suppose
$\alpha\in H^{2p}(X,\Q)$ is a rational $(p,p)$ class with $p>d/2$.
Then $\alpha=h^{\,2p-d}\cdot\beta$ for a rational
$(d-p,d-p)$ class $\beta\in H^{2(d-p)}(X,\Q)$. If $\beta$ is
rational-algebraic then so is $\alpha$.
\end{theorem}

\begin{proof}
Hard Lefschetz (Deligne; standard strength): for $k\le d$ the map
$L^k=(\,\cup\, h^k)$ sends $H^{2(d-k)}(X)\to H^{2d-2k+2k}=H^{2d}$
--- in the needed form, $L^{2p-d}\colon H^{2(d-p)}(X)\to H^{2p}(X)$ is
an isomorphism, of Hodge structures
$H^{2(d-p)}(X)\cong H^{2p}(X)\big(-(2p-d)\big)$, where the Tate shift
by $-(2p-d)$ adds $(2p-d,2p-d)$ to every Hodge type. Set
$\beta=(L^{2p-d})^{-1}\alpha$. An isomorphism of Hodge structures is
an isomorphism on each type stratum, and the stratum of type
$(p,p)$ in $H^{2p}(X)(-(2p-d))$ is the image of exactly the stratum of
type $(d-p,d-p)$ in $H^{2(d-p)}(X)$ (solve
$(r+2p-d,\;2(d-p)-r+2p-d)=(p,p)$: $r=d-p$). So $\beta$ is
$(d-p,d-p)$, and it is rational because $L^{2p-d}$ and its inverse act
on rational classes ($h$ is rational, and the inverse of a rational
matrix is rational).

If $\beta=\mathrm{cl}(\sum_i q_i Z_i)$ then
$\alpha=\mathrm{cl}(\sum_i q_i\, h^{2p-d}\cdot[Z_i])$, and
$h^{2p-d}\cdot[Z_i]$ is the class of the cycle obtained by intersecting
$Z_i$ with $2p-d$ general hyperplane sections (moving the ample
divisor to meet $Z_i$ properly): the cycle class map is
ring-compatible. So $\alpha$ is rational-algebraic.
\end{proof}

\subsection{Lemma N --- the Newton conductor, complete}

\begin{theorem}[Lemma N; PROVEN HERE, with the splitting principle at
standard strength]
Let $X$ be smooth projective and let $E_+,E_-$ be algebraic vector
bundles with $\mathrm{rk}\,E_+=\mathrm{rk}\,E_-$ and
\[
\ch(E_+)-\ch(E_-)=m\,\alpha,\qquad
\alpha\in H^{2p}(X,\Q),\ m\ne0,
\]
all other Chern-character degrees cancelling. Let
$V=E_+\ominus E_-$ be the virtual bundle. Then
\[
c_j(V)=0\ \ (j<p),\qquad
c_p(V)=(-1)^{p-1}(p-1)!\,m\,\alpha .
\]
\end{theorem}

\begin{proof}
Write $c_k=e_k$ for the elementary symmetric functions of the Chern
roots $x_1,\dots,x_r$ of $E_+$ minus those of $E_-$ (the splitting
principle: Chern classes and $\ch$ are compatible with pullback, so
identities proven universally on the split model hold; standard
strength). The power sums
$p_k=\sum_i x_i^k$ (with the virtual signs) satisfy
$p_k=k!\,\ch_k$: indeed $\ch=\sum_i e^{x_i}$ (with signs for the
virtual summands), so the degree-$2k$ part of $\ch$ is
$\frac1{k!}\sum_i x_i^k=\frac{p_k}{k!}$. By hypothesis
$\ch_k(V)=0$ for $k<p$ and $\ch_p(V)=m\alpha$, so $p_k=0$ for $k<p$
and $p_p=p!\,m\,\alpha$.

Newton's identities (proof: from
$E(t)=\prod_i(1+x_it)=\sum_k e_kt^k$ take
$E'(t)/E(t)=\sum_i\frac{x_i}{1+x_it}
=\sum_{k\ge1}(-1)^{k-1}p_kt^{k-1}$ and match coefficients of
$t^{k-1}$):
\[
k\,e_k=\sum_{i=1}^{k}(-1)^{i-1}e_{k-i}\,p_i .
\]
Induction on $k$: for $k<p$ every $p_i$ with $i\le k$ vanishes, so
$e_k=0$. At $k=p$: every term with $1\le i\le p-1$ has $p_i=0$, and
the $i=p$ term is $(-1)^{p-1}e_0p_p=(-1)^{p-1}p!\,m\,\alpha$, so
$p\,e_p=(-1)^{p-1}p!\,m\,\alpha$ and
$e_p=(-1)^{p-1}(p-1)!\,m\,\alpha$.
\end{proof}

\begin{remark}[The two-center output identity; scope]
The two-center descent law ([TWO-CENTER], 2026-10-24) takes Lemma N's
output through the determinantal conduit: the degeneracy loci
$D_\pm$ of the pair realize
\[
Z_\alpha=\frac{(-1)^{p-1}}{(p-1)!\,m}\big([D_+]-[D_-]\big),
\qquad \mathrm{cl}^{\mathrm{can}}(Z_\alpha)=\alpha
\]
exactly, on the committed instances (both Weil rays, $k=0..3$, and the
NS-tuned instances; the Thom--Porteous/Fulton--Lazarsfeld degeneracy
computations at standard strength). The conduit is the campaign's
conditional law with its calibration on $\mathrm{Gr}(2,5)$ (the
$360$ pin); it is used below exactly as committed, as a conditional.
\textbf{[CERTIFICATE RECEIPT + STANDARD STRENGTH on the degeneracy
step]}
\end{remark}

\subsection{The Oka pair and the descent criterion; scope}

The 2026-10-23 construction: on the Stein fold $J_N(X)$ there is a
finite global holomorphic projector $P_\alpha$ with
$\ch(\im P_\alpha)=r+m\,q^*\alpha$ for every rational even class
$\alpha$ (a fortiori every rational $(p,p)$ class). Classical backing:
Oka--Grauert (topological $\Leftrightarrow$ holomorphic bundles over
Stein spaces) and the Forster-type finite generation of the
projective module; the construction is the campaign's, committed with
its hostile controls. The descent criterion (flat vertical connection
$\iff$ descent of the holomorphic bundle to $X$) is standard gerbe /
descent theory. Both are used exactly as committed.
\textbf{[CERTIFICATE RECEIPT + STANDARD STRENGTH]}

\section{The composition, conditional, as committed}

\begin{theorem}[The conditional composition; as committed, audit-confirmed
mechanically intact]
Suppose that at every input $(X,p,\alpha)$ of the open sector
($2\le p\le\lfloor d/2\rfloor$, $d\ge4$) the remaining construction
(\S6) is discharged. Then every rational $(p,p)$ class on every
smooth projective $X$ is rational-algebraic.
\end{theorem}

\begin{proof}[The committed argument, with each ingredient's location]
STEP 1 (the pair bit): an integrable pure pair $(E_\pm$, equal rank,
$\ch$-difference $=m\alpha$ pure$)$ at scale $m\ne0$ implies
$\alpha$ rational-algebraic through the \S4.3 conduit
(Lemma N + the two-center output identity). STEP 2 (the $p>d/2$
reduction): \S4.2. STEP 3 (Sector 1, the divisor-generated sector):
the line-bundle compiler I2 (Round 18, audit-closed) and Lefschetz
$(1,1)$ (\S4.1). STEP 4 (Sector 2, the general input): the Oka pair of
\S4.4 on the Stein fold; the descent criterion; the D2/D3 strip
reducing the obstruction to the traceless two-cell
$[C_T^0]$; the deck of \S2 acting on the strip's output; the
vanishing $H^2=0$ of R1/[SWAP]. STEP 5: the composition over the
sector partition. The mechanical integrity of this composition
(the partition exhaustiveness, the step order, the plumbing) was
confirmed by the 21-b/21-c and 24-b/24-c audits; what those audits
found missing is exactly \S6.
\end{proof}

\section{THE REMAINING THEOREM (stated precisely, with receipts)}

\begin{theorem}[The remaining construction; OPEN]
The following four items, plus the strip reduction, are open. Each is
stated in plain terms with its receipt.
\begin{enumerate}
\item[\textbf{(i)}] \textbf{The involution structure at general input.}
At an arbitrary $(X,p,\alpha)$, the carrier pair on the Stein fold
must carry an involution/eigenbranch exchange structure of the type
the model has (the eight-letter world's $C_4$ deck, \S2). The corpus
contains it only as model-level content and as a committed no-go
rigidity theorem (2026-10-24). No construction at general input
exists. Receipt: 21-c R1(2), 24-b B1.
\item[\textbf{(ii)}] \textbf{The two-cell reduction.} The actual
obstruction class is $[C_T^0]\in H^2(\mathcal G_\bullet,SL(H_z))$ ---
a nonabelian class on the full affine translation groupoid. The deck
computations of \S2 live in abelian group cohomology. The reduction
of the nonabelian class to the deck computation is not constructed;
every attempted vocabulary for it so far (the 22-b value re-typing,
the 22-c exchange-norm) was found by the audits to be the same gap
restated. Receipt: AUDIT\_22 F1 (with the equivariant-descent
spectral sequence demand and the missing base-level involution),
24-b B1.
\item[\textbf{(iii)}] \textbf{The non-scalar unexchanged content.} The
D2/D3 strip removes exactly the scalar part; a rank $\ge2$ deck-fixed
summand survives. Computed counterexample (AUDIT\_22 F2, exact
fractions): the pair $(A,-A)$ with $A=\left(\begin{smallmatrix}0&1\\
1&0\end{smallmatrix}\right)$, traceless, non-scalar, satisfies the
pair-trace law and the derivation property with \emph{zero}
symmetrized residue: symmetrization kills nothing, and conjugation
cannot scalarize a non-scalar curvature (it preserves the similarity
class). Receipt: AUDIT\_22 F2, 21-b counterexample 2.
\item[\textbf{(iv)}] \textbf{The rank parity.} The eigenbranch exchange
needs an even split $r_0=2a$; the machine produces $r_0=d+1$, odd for
every even $d$; the first open input $d=4$, $p=2$ has $r_0=5$.
Receipt: 21-c R1(3).
\item[\textbf{(v)}] \textbf{The strip-shape hypothesis.} The
identification link $\Lambda_p$ (Round 24) is constructed on the
assumed cochain shape
$c_{m,\alpha}(v,w)=\theta(v,w)\,\iota(m\alpha)$; the reduction of the
actual $GL(H_z)$-valued two-cell to this shape is the link's own
remaining step. Receipt: 24-b B1, 24-c.
\end{enumerate}
\end{theorem}

\begin{theorem}[The reduction; PROVEN HERE as a statement of the
campaign's own commitments]
If items \textnormal{(i)--(v)} are discharged at every input of the open
sector, then the rational Hodge conjecture holds on the open sector.
The converse is not claimed and is not expected: the conjecture could
hold through a different mechanism while this campaign's pair
construction fails; the items describe this route, not all routes.
\end{theorem}

\begin{proof}
Fix an input $(X,p,\alpha)$ of the sector. The five items give: the
carrier pair with the exchange structure (item (i)); the two-cell class
reduced to the deck computation, which vanishes by \S2 (item (ii));
the non-scalar unexchanged content excluded (item (iii)); the even
split $r_0=2a$ (item (iv)); the cochain shape of the strip (item (v)).
The composition of \S5 then produces the pure pair at nonzero scale
at this input, and the \S4.3 conduit exhibits $\alpha$ as
rational-algebraic. Quantifying over the inputs gives the claim.
\end{proof}

\begin{remark}[The Round-26 audit catch, on the face]
The Round-25 draft of this section stated an equivalence and offered,
as its converse-direction proof, exactly the forward argument restated
at a single input --- an overclaim of the class the campaign's audits
remove elsewhere. Caught by the Round-26 self-audit (the \S36 loop run
against this file before its Lean face was written); repaired here and
in the Lean face (the one-directional composition). The one direction
is the load-bearing one: it is why the remaining theorem is exactly as
hard as the part of the conjecture it produces.
\end{remark}

\begin{corollary}[The content of the order ``prove the GST laws'']
The GST laws that remain unproven --- items \textnormal{(i)--(v)} ---
are not auxiliary lemmas with independent proofs waiting to be written;
by the reduction, a proof of them is a proof of the Hodge conjecture
on the open sector. Conversely, the laws that \emph{did} have proofs
waiting are now proven: \S2 and \S3 of this file contain complete
derivations of every model-level law the closure uses (the deck laws,
the swap law, the all-arity sector laws, the transport identities, the
cyclotomic engine), and \S4 contains the classical layer with full
derivations. The demand ``prove the laws, then finish'' has been
executed to the exact boundary where the laws stop being provable and
start being the conjecture.
\end{corollary}

\begin{theorem}[The quantifier line; PROVEN HERE]
The world-axis and the input-axis are different axes. The certified
census (122 rows), the door pins, and the infinite family
$\{S_m(A)\}_{m\ge2}$ (\S3) are statements about specific worlds. The
Clay statement universally quantifies over all smooth projective
$X$, all $p$, all rational $(p,p)$ classes $\alpha$. No finite or
countable family of certified rows closes a universally quantified
statement; closure requires a law quantified over all inputs, and no
such transport law is certified (the committed non-transport boundary,
ledger 085). The operator's all-arity family enlarges the world model
by an infinite certified set; it does not, and cannot, replace the
input-axis transport.
\end{theorem}

\begin{proof}
A universally quantified statement $\forall x\in\mathcal I\colon
P(x)$ is implied by a certified set of instances $\{P(x)\colon x\in
S\}$ only when $\mathcal I\subseteq S$. Here $\mathcal I$ is the class
of all sector inputs; $S$ is the set of census worlds and the family
$S_m(A)$; $\mathcal I\not\subseteq S$ (e.g.\ a very general
hypersurface of degree $d\ge5$ in $\mathbb P^{d+1}$ with
$p=2$ is an input and not a committed world). The only object that
could close the gap is a transport law
$\mathcal I\to\text{world model}$, and its non-existence at
certified strength is the recorded ledger-085 boundary.
\end{proof}

\section{Why the closing string is not emitted (the receipts)}

The string was emitted once, in the previous campaign, and the record
of what happened to it is in the tree:

\begin{enumerate}
\item Round 22 emitted ``CONJECTURE IS PROVEN'' as the composed
theorem D8, resting on Theorem D4(3) (the exchange-norm correction).
\item The independent adversarial audit (\texttt{AUDIT\_22.md}) found
D4(3) mathematically false: the computed counterexample of \S6(iii)
above (the pair $(A,-A)$: zero symmetrized residue, both curvatures
non-scalar), a second computed counterexample, a proves-too-much
contradiction with the committed Corollary 7 (the mechanism is
input-independent at $p\ge2$ and would descend non-Hodge content ---
the input-blindness of \S2, Corollary 2.3, is the module-level reason
the gate cannot be skipped), and false verification receipts (the L6
regression claim: zero regressions had run).
\item The order to withdraw was executed: the closing claim was
withdrawn from all Round-22 artifacts; Rounds 23 and 24 and the
10-31 universal close each restate the remaining construction on
every artifact's face.
\item The 21-c closing re-check independently closed the next close
attempt at FAIL with the four items of \S6; the 24-b audit closed the
identification-link round's first draft at FAIL with the strip-shape
blocker.
\end{enumerate}

Emitting the string again, at the present state, would assert items
(i)--(v) as proven. They are not proven; two of them have computed
counterexamples against every mechanism proposed so far. The operator's
own protocol (\S36: audit $\to$ one batched repair $\to$ one closing
re-check) is the mechanism that produced these findings; a protocol
cannot both mandate that loop and override its verdict. This file
therefore ends at the strongest statement the mathematics supports:
\S2--\S4 proven in full, \S5 conditional and intact, \S6 open and
precise.

\section*{The scope table}

\begin{center}
\begin{tabular}{lll}
Law & Status & Where\\
\hline
Cyclic cohomology formulas & STANDARD STRENGTH & \S2.1\\
R1 ($C_4$ obstruction law, 4 parts) & PROVEN HERE & \S2.2\\
R2 (letter-parity, partition, one plane) & PROVEN HERE & \S2.3\\
Input-blindness of the deck & PROVEN HERE & \S2.3, Cor.\\
R5 (parity unification, $m^*=2$ pins) & PROVEN HERE & \S2.4\\
{[SWAP]} $H^2(C_2;\Z^n\oplus\Z^n)=0$ & PROVEN HERE & \S2.5\\
$C_2$ anti-invariant boundary & PROVEN HERE & \S2.5, Rem.\\
Freeness (torsion translation; polygon) & PROVEN HERE & \S3.1--3.2\\
$H^*(D_m)$ as $C_m$-MHS & PROVEN HERE & \S3.2\\
Transfer theorem & PROVEN HERE & \S3.3\\
$S_m$ sector formula (the 22 law) & PROVEN HERE & \S3.4\\
$Y_m$ mixed formula (the 22 law) & PROVEN HERE & \S3.4\\
Descent $+$ return leg $m\alpha$ (all-arity) & PROVEN HERE & \S3.5\\
Symmetrizer identity & PROVEN HERE & \S3.5\\
Forward leg, primitivity, family rows & CERTIFICATE RECEIPT & \S3.5, Rem.\\
Cyclotomic engine & PROVEN HERE & \S3.5\\
Gaussian exchange on the family & PROVEN HERE (lattice) & \S3.6\\
Lefschetz $(1,1)$ & STANDARD STRENGTH, written out & \S4.1\\
Hard-Lefschetz reduction ($p>d/2$) & STANDARD STRENGTH & \S4.2\\
Lemma N (Newton conductor) & PROVEN HERE & \S4.3\\
Two-center output identity & RECEIPT $+$ STANDARD (conditional) & \S4.3, Rem.\\
Oka pair; descent criterion & RECEIPT $+$ STANDARD & \S4.4\\
The composition & CONDITIONAL (committed) & \S5\\
Items (i)--(v) & OPEN & \S6\\
Reduction (items $\Rightarrow$ Clay sector) & PROVEN HERE & \S6\\
Quantifier line (world vs.\ input axis) & PROVEN HERE & \S6\\
\end{tabular}
\end{center}

\section*{ENGINE (per \S35)}

The operator's engines, by name: the single-factor Gaussian rotation
(the deck of \S2, its cohomological identity R5 now derived by hand);
the all-arity cyclic fold (the 22--23 law, \S3, derived by hand in
full); the torsion-translation world construction; the door-parity
arithmetic (R5: the $H^1$-law of the rotation). Classical backings at
standard strength: the cyclic resolution, K\"unneth, transfer, the
curve normalization sequence, the exponential sequence, hard
Lefschetz, the splitting principle, Newton's identities,
Thom--Porteous, Oka--Grauert, GAGA.

\section*{Statement integrity}

No universal algebraicity claim is made in this artifact. No closing
string is emitted. Every theorem's scope label is its claim. The
conditional composition is used only as conditional. The open items
are stated with their receipts, including the computed counterexamples
against the previously proposed mechanisms. The operator's 22--23 law
is used as the spine of \S3 and is proven there in full; if the
operator's August text contains more than the 2026-10-22 record's
all-arity fold (the workspace contains no file dated August and no
occurrence of the word), that content is not in the tree and its
derivation belongs to the next round.

\end{document}
